\chapter{Systematic Literature Review}
\label{ch:Systematic Literature Review}


Kitchenham’s guidelines were originally developed for evidence-based software engineering 
and have become one of the most widely adopted standards for conducting systematic 
literature reviews in computer science and technical research disciplines 
\cite{kitchenham_procedures_2004}. 
The methodology provides detailed, domain-specific procedures for planning, searching, 
selecting, and synthesizing studies in areas where research outputs are often 
heterogeneous—such as algorithmic models, machine learning architectures, or engineering 
pipelines. Unlike biomedical-focused protocols such as PRISMA 
\cite{moher_preferred_2009}, which were designed primarily for clinical trials 
and human-subject research with standardized quantitative outcomes, Kitchenham explicitly 
addresses challenges unique to computational research: diversity of study designs, 
inconsistent terminology, non-uniform evaluation metrics, and the frequent presence of 
code-based or algorithmic artifacts instead of comparable statistical datasets. 
Similarly, frameworks like SALSA \cite{noauthor_typology_nodate} emphasize high-level 
structuring but do not prescribe the step-by-step operational procedures needed for 
reproducibility in technical fields.

For this thesis—situated within deep learning, machine learning, and vector-graphics 
generation—a Kitchenham-style SLR is needed because the literature consists of diverse 
model architectures, varied input modalities, and different evaluation criteria. 
Kitchenham’s protocol supports this complexity by providing structured phases (planning, 
conducting, and reporting) that guide the reviewer through formulating rigorous research 
questions, defining search strategies, operationalizing selection criteria, assessing 
methodological quality, and performing narrative synthesis. This structure minimizes 
bias, increases reproducibility, and ensures that the review captures the full range of 
relevant deep-learning and vectorization approaches. Consequently, Kitchenham offers the 
most appropriate methodological foundation for achieving a systematic, transparent, and 
technically robust synthesis of current research on SVG-related generative models 
\cite{kitchenham_systematic_2009}.


\section{Phase 1 — Planning}

% \subsection{1. Justifying the Need for an SLR}

% To ensure methodological rigor, this study adopts the Kitchenham guidelines for conducting Systematic Literature Reviews (SLRs) in computer science. Several existing works (e.g., \textit{Multiscale Visualization: A Structured Literature Analysis}, \textit{Systematic Review of Synthetic CT Generation Methodologies}, \textit{A Systematic Literature Review of Deep Learning Approaches for Sketch-Based Image Retrieval}) follow similar approaches using Kitchenham or PRISMA-based SLR structures.

% Before conducting the SLR, the following questions must be addressed:
% \begin{itemize}
%    \item Is a similar SLR already published in the domain?
%    \item Is there sufficient primary literature available to support an SLR?
% \end{itemize}

\subsection {Review Protocol}

\paragraph{Databases}

\begin{itemize}
    \item arXiv
    \item Google Scholar
    \item ACM Digital Library
\end{itemize}


\paragraph{Search Strings/Keywords}
\subparagraph{SVG Generation}
\begin{verbatim}
"SVG generation"
"scalable vector graphics" AND generation
"SVG generative models"
"vector graphics generation"
"SVG generation" AND deep learning
"SVG generation" AND machine learning
\end{verbatim}

\subparagraph{Raster-to-Vector Vectorization}
\begin{verbatim}
"raster-to-vector" AND deep learning
"image-to-vector" AND neural network
"image vectorization" AND machine learning
"deep vectorization"
"differentiable vector graphics"
\end{verbatim}

\subparagraph{Text-to-SVG}
\begin{verbatim}
"text-to-SVG"
"text-to-vector graphics"
\end{verbatim}


\subparagraph{Structure of SVG, Editability and Quality}
\begin{verbatim}
"SVG Basics"
"SVG fundamentals"
"Scalable Vector Graphics"
"SVG editability"
"editable SVG"
"SVG path simplification"
"layered SVG" AND generation" 
\end{verbatim}

\subparagraph{Survey and Review Searches}
\begin{verbatim}
"SVG" AND "systematic literature review"
"vector graphics" AND survey
"vector image" AND review
"image vectorization" AND survey
"graphic generation" AND review
"SVG Generation" AND benchmark
"image vectorization" AND benchmark
\end{verbatim}



% \paragraph{Population (P):}
% \begin{itemize}
%     \item SVG properties
%     \item SVG generation methods
%     \item Deep learning and machine learning architectures used for SVG generation
%     \item Output structures
%     \item Evaluation approaches
%     \item Quality across models
% \end{itemize}

% \paragraph{Intervention (I):}
% \begin{itemize}
%     \item Diffusion-based or generative models that output SVG-like structures
%     \item Sketch-to-SVG pipelines
%     \item Image-to-SVG vectorization pipelines
%     \item LLM-based SVG generation
%     \item Hybrid multimodal SVG generation (text + image)
%     \item ML and DL architectures (Transformers, VAEs, autoregressive models)
% \end{itemize}

% \paragraph{Comparison (C):}
% \begin{itemize}
%     \item Comparison across different deep learning models
%     \item Rule-based vectorization vs. ML-based vectorization
%     \item Structured SVG generation vs. raster-only generative models
%     \item Different input modalities (sketch, raster image, text, multimodal)
%     \item Editability of SVG outputs
% \end{itemize}

% \paragraph{Outcomes (O):}

% Primary outcomes:
% \begin{itemize}
%     \item SVG quality
%     \item SVG editability (e.g., number of nodes, clean paths, human-editable structure)
%     \item Structural characteristics (hierarchy, grouping, path consistency)
%     \item Reported performance metrics
%     \item Comparative performance between open-source models
% \end{itemize}

% Secondary outcomes:
% \begin{itemize}
%     \item Practical usability for downstream editing or integration
% \end{itemize}

% \subsection{3. Creating the Review Protocol}

% \paragraph{3.1 Search Strategy}

% \textbf{Databases:}
% \begin{itemize}
%     \item ACM Digital Library
%     \item IEEE Xplore
%     \item arXiv
%     \item Google Scholar
% \end{itemize}

% \textbf{Example Search Strings:}
% \begin{verbatim}
% ("SVG" OR "vector representation" OR "vectorized image" OR
%  "vector data representation" OR "Scalable Vector Graphics" OR
%  "vectorization" OR "image-to-vector" OR "raster-to-vector" OR
%  "raster-to-SVG" OR "vector structure analysis" OR "sketch-to-vector" OR
%  "text-to-SVG" OR "SVG generation" OR
%  "deep learning" OR "machine learning" OR "transformer" OR "diffusion" OR
%  "generative model" OR "vector quality evaluation" OR
%  "survey" OR "systematic literature review")
% \end{verbatim}

% Search fields include title, abstract, and keywords.

\paragraph{Study Selection Criteria}

\subparagraph{Inclusion Criteria:}
\begin{itemize}
    \item \textbf{Language and Publication:} Written in English and published in a peer-reviewed journal, conference, or workshop.
    \item \textbf{Time Range:} Published between 2020--2025 (post-emergence of modern SVG generative models such as DeepSVG).
    \item \textbf{Domain Relevance:} Falls within computer vision, computer graphics, AI, machine learning, or deep learning.
    \item \textbf{Reproducibility:} Provides an open-source implementation or references a publicly available repository. Availability of pre-trained weights is preferred, but source code alone is sufficient if the method is reproducible.
    \item \textbf{SVG Output Requirement:} Describes a model, method, or architecture that directly generates SVG or structured vector outputs, including:
        \begin{itemize}
            \item Bézier curves,
            \item SVG \texttt{<path>} commands,
            \item layered or hierarchical SVG structure,
            \item vector strokes or SVG-like stroke sequences,
            \item structured outputs suitable for editability analysis.
        \end{itemize}
    \item \textbf{Input Modalities:} Supports at least one of:
        \begin{itemize}
            \item image-to-SVG vectorization,
            \item text-to-SVG generation,
            \item multimodal SVG generation,
        \end{itemize}
    \item \textbf{Architectural Transparency:} Provides sufficient technical detail to extract the model architecture or internal SVG representation (e.g., decoder design, transformer layers, vector tokenization, differentiable rasterization, or geometric parameterization).
    \item \textbf{Evaluation:} Provides quantitative metrics and/or qualitative SVG examples that allow analysis of structural quality and editability.
\end{itemize}

\subparagraph{Exclusion Criteria:}
\begin{itemize}
    \item \textbf{No Vector Output:} The method does not generate vector or SVG outputs (pixel-only image generation).
    \item \textbf{Non-ML Methods:} Purely rule-based or classical algorithmic vectorization methods without machine learning or deep learning.
    \item \textbf{Closed-Source Systems:} No publicly available code, weights, or sufficient technical detail for reproduction.
    \item \textbf{Insufficient Evaluation:} Papers lacking SVG evaluation, qualitative vector examples, or any structural analysis of outputs.
    \item \textbf{Out-of-Scope Vector Work:} Studies focused solely on vector editing, rendering, compression, markup optimization, or stylistic modification without SVG generation.
    \item \textbf{Domain Mismatch:} 
        \begin{itemize}
            \item handwriting synthesis,
            \item font generation,
            \item datasets,
            \item sketch generation not intended for SVG graphics.
        \end{itemize}
    \item \textbf{Temporal Exclusion:} Studies published before 2020, except when cited for historical context only.
    \item \textbf{LLM Restriction:} Proprietary LLMs (e.g., GPT-4, Gemini, Claude) that can emit SVG text but do not provide open-source architecture, reproducible training details, or transparent vector representations.
\end{itemize}


\paragraph{Study Selection Procedure}

\begin{itemize}
    \item Initial records
    \item Duplicate removal
    \item Title/abstract screening
    \item Full-text assessment
    \item Final included studies
\end{itemize}

\paragraph{Quality Assessment Criteria}

Based on Kitchenham:
\begin{itemize}
    \item Clear aims and objectives
    \item Clear description of methodology
    \item Clear model/method architecture
    \item Appropriate study design
    \item Evaluation methodology
    \item Clear results
    \item Discussion of limitations
    \item Open-source implementation
\end{itemize}

\paragraph{Data Extraction Plan}

\begin{itemize}
    \item Bibliographic metadata
    \item Study context and goals
    \item Model/method characteristics
    \item Vector structure and SVG properties
    \item Evaluation metrics and results
    \item Reproducibility details
\end{itemize}

% \paragraph{Data Synthesis Plan}

% Narrative synthesis includes:
% \begin{itemize}
%     \item Categorization by input modality, model type, output structure, evaluation type, reproducibility
%     \item Identification of evaluation gaps
%     \item Analysis of structural/editability challenges
%     \item Architectural pattern comparison
%     \item Model selection analysis (suitability, runtime feasibility)
% \end{itemize}

% \paragraph{Protocol Review}

% The protocol may be reviewed by supervisors prior to execution.

\section{Phase 2 — Conducting the Review}

\subsection{Execution of the Search Strategy}
The search strings defined in Phase~1 were executed across the selected databases 
(Google Scholar, arXiv, and ACM Digital Library) to identify relevant literature on 
SVG generation and vector-based deep learning methods. Searches were performed using 
the conceptual string groups defined in the search strategy, with database-specific 
syntax adjustments where necessary. All records retrieved by these searches were 
exported and documented before any filtering or removal of duplicates. Table~\ref{tab:searchlog} 
summarises the number of studies returned by each search group and database.

\begin{table}[h!]
\centering
\caption{Search Log Summary Across Databases and Conceptual Groups}
\label{tab:searchlog}
\begin{tabular}{lcccc}
\hline
\textbf{Search Group} & \textbf{Google Scholar} & \textbf{arXiv} & \textbf{ACM DL} & \textbf{Total} \\
\hline
SVG Generation & 18 & 15 & 4 & 37 \\
Raster-to-Vector & 10 & 8 & 2 & 20 \\
Text-to-SVG / Multimodal & 7 & 10 & 5 & 22 \\
Surveys/Benchmarks/Reviews & 6 & 3 & 1 & 10 \\
SVG Fundamentals / Background & 3 & 3 & 2 & 8 \\
\hline
\textbf{Total Hits} & \textbf{44} & \textbf{39} & \textbf{12} & \textbf{95} \\
\hline
\end{tabular}
\end{table}

\subsection{Removing of Duplications}
After merging all search results, duplicate records—including repeated titles and preprint–conference duplicates—were removed. This resulted in 72 unique studies.

\subsection{5. Study Selection}
\begin{itemize}
    \item Stage 1: Title, Keywords and Abstract screening.

After removing duplicates, \textbf{72} unique studies remained. A brief title and 
abstract screening excluded works clearly unrelated to SVG generation (e.g., 
sketch-only models, raster methods, classical vectorization, or closed-source LLMs), 
resulting in \textbf{38} papers proceeding to full-text review.


    \item Stage 2: Full-text screening
\end{itemize}

\subsection{6. Study Quality Assessment}
Apply the quality checklist.

\subsection{7. Data Extraction}
Use standardized extraction tables.

\subsection{8. Data Synthesis}

\paragraph{8.1 Narrative Synthesis}
Summarize findings, patterns, and categories.

\paragraph{8.2 Sensitivity Analysis}
Remove low-quality studies and reassess.

\section{Phase 3 — Reporting}

\subsection{9. Prepare the SLR Report}
Document all methodology, results, synthesis, and identified gaps.
